% !TeX root = part01.tex

\documentclass[../final.tex]{subfiles}

\begin{document}

% ============================================================
% Практическое занятие 1
% ============================================================

\practice{1}{03.09.2026}{Кинематика}


% ============================================================
% Релятивистская кинематика
% ============================================================

\section*{Релятивистская кинематика}

Для релятивистской частицы выпишем основные соотношения между
энергией, импульсом, массой и скоростью:
%
\begin{equation*}
    E=\gamma mc^2,
\end{equation*}
%
где
%
\begin{equation*}
    \gamma
    =
    \frac{1}{\sqrt{1-\beta^2}},
    \qquad
    \beta
    =
    \frac{v}{c}.
\end{equation*}

Энергия и импульс связаны соотношением
%
\begin{equation*}
    E^2
    =
    (pc)^2+(mc^2)^2.
\end{equation*}

В дальнейшем удобно использовать четырёхмерный импульс частицы.
В системе единиц $c=1$ он имеет вид
%
\begin{equation*}
    P
    =
    (E,\vec p).
\end{equation*}

Его квадрат равен
%
\begin{equation*}
    P^2
    =
    E^2-p_x^2-p_y^2-p_z^2
    =
    m^2.
\end{equation*}

Величина
%
\begin{equation*}
    P^2=m^2
\end{equation*}
%
является \rtxt{3}{лоренц-инвариантом}, то есть не изменяется
при переходе из одной инерциальной системы отсчёта в другую.


% ------------------------------------------------------------
% Система из двух частиц
% ------------------------------------------------------------

\subsection*{Система из двух частиц}

Пусть имеются две частицы с четырёхимпульсами
%
\begin{equation*}
    P_1
    =
    (E_1,\vec p_1),
    \qquad
    P_2
    =
    (E_2,\vec p_2).
\end{equation*}

Полный четырёхимпульс системы равен
%
\begin{equation*}
    P_1+P_2
    =
    (E_1+E_2,\vec p_1+\vec p_2).
\end{equation*}

Рассмотрим его квадрат:
%
\begin{align*}
    (P_1+P_2)^2
    &=
    (E_1+E_2)^2
    -
    (\vec p_1+\vec p_2)^2
    \\[3pt]
    &=
    E_1^2
    +
    E_2^2
    +
    2E_1E_2
    -
    p_1^2
    -
    p_2^2
    -
    2\vec p_1\cdot\vec p_2.
\end{align*}

Так как
%
\begin{equation*}
    E_1^2-p_1^2=m_1^2,
    \qquad
    E_2^2-p_2^2=m_2^2,
\end{equation*}
%
а
%
\begin{equation*}
    \vec p_1\cdot\vec p_2
    =
    p_1p_2\cos\theta,
\end{equation*}
%
то
%
\begin{equation*}
    (P_1+P_2)^2
    =
    m_1^2+m_2^2
    +
    2E_1E_2
    -
    2p_1p_2\cos\theta.
\end{equation*}

Этот же инвариант можно записать через инвариантную массу системы
двух частиц:
%
\begin{equation*}
    (P_1+P_2)^2
    =
    P_{12}^2
    =
    m_{12}^2.
\end{equation*}

Здесь $\theta$ --- угол между импульсами $\vec p_1$ и $\vec p_2$.


% ============================================================
% Преобразования Лоренца
% ============================================================

\section*{Преобразования Лоренца}

Рассмотрим две системы отсчёта, одна из которых движется относительно
другой вдоль оси $x$.

Преобразование Лоренца имеет вид
%
\begin{equation*}
    \begin{pmatrix}
        ct\\
        x\\
        y\\
        z
    \end{pmatrix}
    =
    \begin{pmatrix}
        \gamma      & \beta\gamma & 0 & 0\\
        \beta\gamma & \gamma      & 0 & 0\\
        0           & 0           & 1 & 0\\
        0           & 0           & 0 & 1
    \end{pmatrix}
    \begin{pmatrix}
        ct'\\
        x'\\
        y'\\
        z'
    \end{pmatrix}.
\end{equation*}

Нештрихованные величины относятся к
\rtxt{3}{лабораторной системе отсчёта},
а штрихованные --- к
\rtxt{3}{движущейся системе отсчёта}.

Та же матрица преобразования будет использоваться далее и для
четырёхимпульса.


% ============================================================
% Задача 1
% ============================================================

\newpage
\section{Задача 1. Двухчастичный распад}

Рассмотрим распад частицы массы $M$ на две частицы масс $m_1$ и $m_2$.

\begin{rbox}[left=18pt,right=18pt,top=14pt,bottom=10pt]{[]}
    \centering

    \begin{tikzpicture}[
        >=Latex,
        every node/.style={text=rtext}
    ]
        \coordinate (O) at (0,0);

        \fill[rc3] (O) circle (2.2pt);

        \node[above=5pt] at (O) {$M$};

        \draw[
            ->,
            draw=rtext,
            line width=0.9pt
        ]
        (O) -- (-3,0);

        \draw[
            ->,
            draw=rtext,
            line width=0.9pt
        ]
        (O) -- (3,0);

        \node[above] at (-1.7,0) {$P_1$};
        \node[above] at (1.7,0) {$P_2$};

        \node[below] at (-2.7,0) {$m_1$};
        \node[below] at (2.7,0) {$m_2$};
    \end{tikzpicture}

    \captionsetup{
        font=footnotesize,
        labelfont=bf,
        skip=5pt
    }

    \captionof{figure}{
        Двухчастичный распад частицы массы $M$
        в её системе покоя.
    }
    \label{fig:practice02_decay}
\end{rbox}

Пусть исходная частица покоится. Закон сохранения четырёхимпульса:
%
\begin{equation*}
    P_0
    =
    P_1+P_2.
\end{equation*}

Отсюда
%
\begin{equation*}
    P_2
    =
    P_0-P_1.
\end{equation*}

В системе покоя исходной частицы
%
\begin{equation*}
    P_0
    =
    (M,\vec 0),
\end{equation*}
%
а для продуктов распада
%
\begin{equation*}
    P_1
    =
    (E_1,\vec p_1),
    \qquad
    P_2
    =
    (E_2,\vec p_2).
\end{equation*}

Поэтому
%
\begin{equation*}
    (M,\vec 0)
    =
    (E_1,\vec p_1)
    +
    (E_2,\vec p_2).
\end{equation*}

Чтобы найти энергию первой частицы, воспользуемся равенством
%
\begin{equation*}
    P_2=P_0-P_1
\end{equation*}
%
и возведём обе части в квадрат:
%
\begin{equation*}
    P_2^2
    =
    (P_0-P_1)^2.
\end{equation*}

Так как
%
\begin{equation*}
    P_2^2=m_2^2,
    \qquad
    P_0^2=M^2,
    \qquad
    P_1^2=m_1^2,
\end{equation*}
%
получаем
%
\begin{equation*}
    m_2^2
    =
    M^2+m_1^2-2P_0P_1.
\end{equation*}

В системе покоя исходной частицы
%
\begin{equation*}
    P_0P_1
    =
    ME_1,
\end{equation*}
%
поэтому
%
\begin{equation*}
    m_2^2
    =
    M^2+m_1^2-2ME_1.
\end{equation*}

Отсюда
%
\begin{equation*}
    2ME_1
    =
    M^2+m_1^2-m_2^2,
\end{equation*}
%
и
%
\begin{equation*}
    E_1
    =
    \frac{M^2+m_1^2-m_2^2}{2M}.
\end{equation*}

Аналогично для второй частицы:
%
\begin{equation*}
    E_2
    =
    \frac{M^2+m_2^2-m_1^2}{2M}.
\end{equation*}

Импульсы продуктов распада находятся из релятивистского соотношения
%
\begin{equation*}
    E_i^2
    =
    p_i^2+m_i^2.
\end{equation*}

Следовательно,
%
\begin{equation*}
    p_{1,2}
    =
    \sqrt{
        E_{1,2}^2-m_{1,2}^2
    }.
\end{equation*}

Поскольку исходная частица покоилась,
%
\begin{equation*}
    \vec p_1+\vec p_2=0,
\end{equation*}
%
поэтому
%
\begin{equation*}
    \vec p_1=-\vec p_2,
\end{equation*}
%
и модули импульсов совпадают:
%
\begin{equation*}
    p_1=p_2.
\end{equation*}


% ============================================================
% Задача 2
% ============================================================

\newpage
\section{Задача 2. Неподвижная мишень}

Рассмотрим столкновение налетающей частицы массы $m_1$
с покоящейся частицей массы $m_2$, в результате которого
образуется система массы $M$.

\begin{rbox}[left=18pt,right=18pt,top=14pt,bottom=10pt]{[]}
    \centering

    \begin{tikzpicture}[
        >=Latex,
        every node/.style={text=rtext}
    ]
        \fill[rc3] (-2.4,0) circle (2.2pt);
        \fill[rtext] (0.7,0) circle (2.2pt);

        \draw[
            ->,
            draw=rtext,
            line width=0.9pt
        ]
        (-2.4,0) -- (-0.3,0);

        \draw[
            ->,
            draw=rtext,
            line width=0.9pt
        ]
        (0.7,0) -- (3.0,0);

        \node[above] at (-2.4,0) {$m_1$};
        \node[below] at (-1.4,0) {$E_1,\vec p_1$};

        \node[above] at (0.7,0) {$m_2$};
        \node[above] at (2.5,0) {$M$};
    \end{tikzpicture}

    \captionsetup{
        font=footnotesize,
        labelfont=bf,
        skip=5pt
    }

    \captionof{figure}{
        Столкновение налетающей частицы
        с неподвижной мишенью.
    }
    \label{fig:practice02_fixed_target}
\end{rbox}

Закон сохранения четырёхимпульса:
%
\begin{equation*}
    P_1+P_2=P.
\end{equation*}

В компонентах:
%
\begin{equation*}
    (E_1,\vec p_1)
    +
    (E_2,\vec 0)
    =
    (E,\vec p).
\end{equation*}

Для пространственных компонент
%
\begin{equation*}
    \vec p_1+\vec 0
    =
    \vec p.
\end{equation*}

Так как вторая частица покоится,
%
\begin{equation*}
    E_2=m_2.
\end{equation*}

Возведём закон сохранения четырёхимпульса в квадрат:
%
\begin{equation*}
    (P_1+P_2)^2
    =
    P^2.
\end{equation*}

Правая часть равна
%
\begin{equation*}
    P^2=M^2.
\end{equation*}

Для левой части
%
\begin{equation*}
    (P_1+P_2)^2
    =
    P_1^2+P_2^2+2P_1P_2.
\end{equation*}

Так как
%
\begin{equation*}
    P_1^2=m_1^2,
    \qquad
    P_2^2=m_2^2,
\end{equation*}
%
а покоящаяся вторая частица имеет четырёхимпульс
%
\begin{equation*}
    P_2=(m_2,\vec 0),
\end{equation*}
%
то
%
\begin{equation*}
    P_1P_2
    =
    E_1m_2.
\end{equation*}

Следовательно,
%
\begin{equation*}
    m_1^2+m_2^2+2E_1m_2
    =
    M^2.
\end{equation*}

Отсюда
%
\begin{equation*}
    2E_1m_2
    =
    M^2-m_1^2-m_2^2,
\end{equation*}
%
поэтому
%
\begin{equation*}
    E_1
    =
    \frac{
        M^2-m_1^2-m_2^2
    }{
        2m_2
    }.
\end{equation*}


% ------------------------------------------------------------
% Частный случай: J/psi
% ------------------------------------------------------------

\subsection*{Частный случай: рождение $J/\psi$}

Для частицы
%
\begin{equation*}
    J/\psi,
    \qquad
    M_{J/\psi}
    \sim
    3\,\text{ГэВ},
\end{equation*}
%
рассмотрим столкновение электрона и позитрона:
%
\begin{equation*}
    m_1=m_2=m_e,
    \qquad
    m_e\ll M_{J/\psi}.
\end{equation*}

В случае неподвижной мишени предыдущая формула даёт
%
\begin{equation*}
    E_1
    =
    \frac{
        M_{J/\psi}^2-2m_e^2
    }{
        2m_e
    }.
\end{equation*}

Так как
%
\begin{equation*}
    m_e\ll M_{J/\psi},
\end{equation*}
%
можно пренебречь членом $2m_e^2$:
%
\begin{equation*}
    E_1
    \approx
    \frac{
        M_{J/\psi}^2
    }{
        2m_e
    }.
\end{equation*}

Теперь рассмотрим \rtxt{3}{встречные пучки}.
Если энергии частиц одинаковы, то для рождения частицы массы $M$
энергия каждой частицы должна быть
%
\begin{equation*}
    E_1'
    =
    \frac{M}{2}.
\end{equation*}

Следовательно,
%
\begin{equation*}
    E_1'
    =
    \frac{
        M_{J/\psi}
    }{2}.
\end{equation*}

Сравним требуемые энергии:
%
\begin{align*}
    \frac{E_1}{E_1'}
    &=
    \frac{
        \dfrac{M_{J/\psi}^2}{2m_e}
    }{
        \dfrac{M_{J/\psi}}{2}
    }
    \\[4pt]
    &=
    \frac{
        M_{J/\psi}
    }{
        m_e
    }
    \\[4pt]
    &\approx
    6000.
\end{align*}

Таким образом, при использовании встречных пучков требуемая энергия
одной частицы оказывается примерно в $6000$ раз меньше, чем
в эксперименте с неподвижной мишенью.


% ============================================================
% Задача 3
% ============================================================

\newpage
\section{Задача 3. Открытие антипротона. BEVATRON}

Рассмотрим порог рождения антипротона при столкновении протона
с неподвижным протоном.

\begin{rbox}[left=18pt,right=18pt,top=14pt,bottom=10pt]{[]}
    \centering

    \begin{tikzpicture}[
        >=Latex,
        every node/.style={text=rtext}
    ]
        \fill[rc3] (-2.5,0) circle (2.2pt);
        \fill[rtext] (0.5,0) circle (2.2pt);

        \draw[
            ->,
            draw=rtext,
            line width=0.9pt
        ]
        (-2.5,0) -- (-0.2,0);

        \node[above] at (-2.5,0) {$p$};
        \node[below] at (-1.5,0) {$E,\vec p$};

        \node[above] at (0.5,0) {$p$};
        \node[below] at (0.5,0) {мишень};

        \node[right] at (1.5,0) {
            $p+p\rightarrow \bar p+p+p+p$
        };
    \end{tikzpicture}

    \captionsetup{
        font=footnotesize,
        labelfont=bf,
        skip=5pt
    }

    \captionof{figure}{
        Пороговое рождение антипротона
        на неподвижной протонной мишени.
    }
    \label{fig:practice02_antiproton}
\end{rbox}

Минимальный процесс рождения антипротона имеет вид
%
\begin{equation*}
    p+p
    \rightarrow
    \bar p+p+p+p.
\end{equation*}

В конечном состоянии должны находиться три протона и один антипротон.

В лабораторной системе отсчёта начальный четырёхимпульс:
%
\begin{equation*}
    (E,\vec p)
    +
    (m_p,\vec 0).
\end{equation*}

На пороге рождения в системе центра масс конечные частицы покоятся
относительно друг друга. Поэтому минимальная инвариантная масса
конечного состояния равна
%
\begin{equation*}
    M_{\text{кон}}
    =
    4m_p.
\end{equation*}

Для определения пороговой энергии воспользуемся инвариантностью
квадрата полного четырёхимпульса:
%
\begin{equation*}
    s
    =
    \left[
        (E,\vec p)
        +
        (m_p,\vec 0)
    \right]^2.
\end{equation*}

Раскрывая квадрат, получаем
%
\begin{align*}
    s
    &=
    (E,\vec p)^2
    +
    (m_p,\vec 0)^2
    +
    2(E,\vec p)(m_p,\vec 0)
    \\[3pt]
    &=
    m_p^2
    +
    m_p^2
    +
    2Em_p
    \\[3pt]
    &=
    2m_p^2+2Em_p.
\end{align*}

На пороге рождения
%
\begin{equation*}
    s
    =
    (4m_p)^2
    =
    16m_p^2.
\end{equation*}

Следовательно,
%
\begin{equation*}
    2m_p^2+2Em_p
    =
    16m_p^2.
\end{equation*}

Отсюда
%
\begin{equation*}
    2Em_p
    =
    14m_p^2,
\end{equation*}
%
и
%
\begin{equation*}
    E
    =
    7m_p
    \approx
    7\,\text{ГэВ}.
\end{equation*}

\textbf{Комментарий.}
В исходной записи в промежуточной строке было записано
%
\begin{equation*}
    2m_p^2+2Em_p
    =
    14m_p^2,
\end{equation*}
%
однако следующий результат $E=7m_p$ требует
%
\begin{equation*}
    2m_p^2+2Em_p
    =
    16m_p^2.
\end{equation*}

Число $14$ возникает уже после переноса двух начальных
$m_p^2$ в правую часть.


% ============================================================
% Задача 4
% ============================================================

\newpage
\section{Задача 4. Релятивистский мюон}

Рассмотрим мюон с энергией
%
\begin{equation*}
    E
    =
    1\,\text{ГэВ}.
\end{equation*}

Требуется определить его время жизни $\tau$ в лабораторной системе
отсчёта и расстояние $x$, которое он успеет пролететь.

\begin{rbox}[left=18pt,right=18pt,top=14pt,bottom=10pt]{[]}
    \centering

    \begin{tikzpicture}[
        >=Latex,
        every node/.style={text=rtext}
    ]
        \fill[rc3] (-2.3,0) circle (2.3pt);

        \draw[
            ->,
            draw=rtext,
            line width=0.9pt
        ]
        (-2.3,0) -- (2.8,0);

        \node[above] at (-2.3,0) {$\mu^-$};
        \node[below] at (0.2,0) {$E=1\,\text{ГэВ}$};
    \end{tikzpicture}

    \captionsetup{
        font=footnotesize,
        labelfont=bf,
        skip=5pt
    }

    \captionof{figure}{
        Релятивистский мюон в лабораторной системе отсчёта.
    }
    \label{fig:practice02_muon}
\end{rbox}

Масса мюона приблизительно равна
%
\begin{equation*}
    m_\mu
    \approx
    100\,\text{МэВ}.
\end{equation*}

Собственное время жизни мюона:
%
\begin{equation*}
    \tau'
    =
    2.2\cdot10^{-6}\,\text{с}.
\end{equation*}

Определим лоренц-фактор:
%
\begin{equation*}
    E
    =
    \gamma m_\mu.
\end{equation*}

Следовательно,
%
\begin{equation*}
    \gamma
    =
    \frac{E}{m_\mu}.
\end{equation*}

При
%
\begin{equation*}
    E
    =
    1\,\text{ГэВ}
    =
    1000\,\text{МэВ}
\end{equation*}
%
и
%
\begin{equation*}
    m_\mu
    \approx
    100\,\text{МэВ}
\end{equation*}
%
получаем
%
\begin{equation*}
    \gamma
    \approx
    10.
\end{equation*}

Из-за релятивистского замедления времени
%
\begin{equation*}
    c\tau
    =
    \gamma c\tau'.
\end{equation*}

Следовательно,
%
\begin{equation*}
    \tau
    =
    \gamma\tau'.
\end{equation*}

Подставляя $\gamma=10$,
%
\begin{align*}
    \tau
    &=
    10\cdot2.2\cdot10^{-6}\,\text{с}
    \\[3pt]
    &=
    22\cdot10^{-6}\,\text{с}
    \\[3pt]
    &=
    22\,\text{мкс}.
\end{align*}

Теперь найдём расстояние, которое пролетит мюон:
%
\begin{equation*}
    x
    =
    v\tau.
\end{equation*}

Поскольку
%
\begin{equation*}
    v=\beta c,
\end{equation*}
%
то
%
\begin{equation*}
    x
    =
    \beta c\tau.
\end{equation*}

Связь между $\beta$ и $\gamma$:
%
\begin{equation*}
    \gamma
    =
    \frac{1}{\sqrt{1-\beta^2}}.
\end{equation*}

Отсюда
%
\begin{equation*}
    \beta
    =
    \sqrt{
        1-\frac{1}{\gamma^2}
    }.
\end{equation*}

Поэтому
%
\begin{equation*}
    x
    =
    \sqrt{
        1-\frac{1}{\gamma^2}
    }\,
    c\tau.
\end{equation*}

Так как
%
\begin{equation*}
    \tau
    =
    \gamma\tau',
\end{equation*}
%
имеем
%
\begin{equation*}
    x
    =
    \sqrt{
        1-\frac{1}{\gamma^2}
    }\,
    \gamma c\tau'.
\end{equation*}

При $\gamma=10$
%
\begin{equation*}
    x
    \approx
    10c\tau'.
\end{equation*}

Подставляя
%
\begin{equation*}
    c
    \approx
    3\cdot10^8\,
    \frac{\text{м}}{\text{с}},
    \qquad
    \tau'
    =
    2.2\cdot10^{-6}\,\text{с},
\end{equation*}
%
получаем величину порядка нескольких километров:
%
\begin{equation*}
    x
    \approx
    6\,\text{км}.
\end{equation*}


% ============================================================
% Задача 5
% ============================================================

\newpage
\section{Задача 5. Преобразование углов}

Рассмотрим частицу, которая в движущейся системе отсчёта имеет
импульс $\vec p'$, направленный под углом $\theta'$ к оси $x'$.
Сама система $x'y'$ движется относительно лабораторной системы
вдоль оси $x$ со скоростью $v$.

\begin{rbox}[left=18pt,right=18pt,top=14pt,bottom=10pt]{[]}
    \centering

    \begin{tikzpicture}[
        >=Latex,
        every node/.style={text=rtext}
    ]
        % =====================================================
        % ДО СТОЛКНОВЕНИЯ
        % =====================================================

        % Налетающая частица M
        \fill[rc3] (-4.2,0) circle (2.2pt);
        \node[above left=1pt] at (-4.2,0) {$M$};

        \draw[
            ->,
            draw=rtext,
            line width=0.9pt
        ]
        (-4.15,0) -- (-2.6,0);

        \node[below] at (-3.4,0) {$\vec p_0$};

        % Мишень m
        \fill[rtext] (-1.8,0) circle (2.2pt);
        \node[above right=1pt] at (-1.8,0) {$m$};

        % Стрелка перехода к конечному состоянию
        \draw[
            ->,
            draw=rc3,
            line width=0.8pt
        ]
        (-1.0,0) -- (0.3,0);

        % =====================================================
        % ПОСЛЕ СТОЛКНОВЕНИЯ
        % =====================================================

        % Точка разлёта
        \coordinate (O) at (1.3,0);

        % Горизонтальная опорная линия для угла
        \draw[
            draw=rtext,
            line width=0.7pt
        ]
        (O) -- ++(1.2,0);

        % Верхняя частица m с импульсом P2
        \draw[
            ->,
            draw=rtext,
            line width=0.9pt
        ]
        (O) -- ++(2.6,1.8);

        \node[above left=2pt] at (2.25,0.72) {$m$};
        \node[above right=2pt] at (3.75,1.63) {$\vec P_2$};

        % Нижняя частица M с импульсом P1
        \draw[
            ->,
            draw=rtext,
            line width=0.9pt
        ]
        (O) -- ++(2.9,-1.15);

        \node[below left=1pt] at (2.35,-0.42) {$M$};
        \node[below right=2pt] at (3.95,-1.02) {$\vec P_1$};

        % Угол theta между горизонтальной опорой и P2
        \draw[
            draw=rtext,
            line width=0.8pt
        ]
        ($(O)+(0.75,0)$)
        arc[
            start angle=0,
            end angle=35,
            radius=0.75
        ];

        \node at ($(O)+(0.98,0.34)$) {$\theta$};

    \end{tikzpicture}

    \captionsetup{
        font=footnotesize,
        labelfont=bf,
        skip=5pt
    }

    \captionof{figure}{
        Упругое рассеяние частицы массы $M$
        на покоящейся частице массы $m$.
    }
    \label{fig:practice02_scattering}
\end{rbox}

Нужно определить угол $\theta$ в лабораторной системе отсчёта.

Четырёхимпульс частицы в движущейся системе имеет компоненты
%
\begin{equation*}
    \begin{pmatrix}
        E'/c\\
        p_x'\\
        p_y'\\
        p_z'
    \end{pmatrix}
    =
    \begin{pmatrix}
        E'/c\\
        p'\cos\theta'\\
        p'\sin\theta'\\
        0
    \end{pmatrix}.
\end{equation*}

Переход в лабораторную систему выполняется преобразованием Лоренца:
%
\begin{equation*}
    \begin{pmatrix}
        E/c\\
        p_x\\
        p_y\\
        p_z
    \end{pmatrix}
    =
    \begin{pmatrix}
        \gamma      & \beta\gamma & 0 & 0\\
        \beta\gamma & \gamma      & 0 & 0\\
        0           & 0           & 1 & 0\\
        0           & 0           & 0 & 1
    \end{pmatrix}
    \begin{pmatrix}
        E'/c\\
        p'\cos\theta'\\
        p'\sin\theta'\\
        0
    \end{pmatrix}.
\end{equation*}

В движущейся системе угол определяется как
%
\begin{equation*}
    \tan\theta'
    =
    \frac{p_y'}{p_x'}.
\end{equation*}

В лабораторной системе
%
\begin{equation*}
    \tan\theta
    =
    \frac{p_y}{p_x}.
\end{equation*}

Из преобразования Лоренца для продольной компоненты получаем
%
\begin{equation*}
    p_x
    =
    \beta\gamma\frac{E'}{c}
    +
    \gamma p'\cos\theta'.
\end{equation*}

Поперечная компонента не изменяется:
%
\begin{equation*}
    p_y
    =
    p'\sin\theta'.
\end{equation*}

Следовательно,
%
\begin{equation*}
    \tan\theta
    =
    \frac{
        p'\sin\theta'
    }{
        \beta\gamma\dfrac{E'}{c}
        +
        \gamma p'\cos\theta'
    }.
\end{equation*}

Вынесем $\gamma p'$ из знаменателя:
%
\begin{equation*}
    \tan\theta
    =
    \frac{
        p'\sin\theta'
    }{
        \gamma p'
        \left(
            \beta\frac{E'}{cp'}
            +
            \cos\theta'
        \right)
    }.
\end{equation*}

Сокращая $p'$, получаем
%
\begin{equation*}
    \tan\theta
    =
    \frac{
        \sin\theta'
    }{
        \gamma
        \left(
            \beta\frac{E'}{cp'}
            +
            \cos\theta'
        \right)
    }.
\end{equation*}

Теперь воспользуемся соотношениями
%
\begin{equation*}
    E'
    =
    \gamma' mc^2,
\end{equation*}
%
и
%
\begin{equation*}
    p'
    =
    \gamma' mv'.
\end{equation*}

Тогда
%
\begin{align*}
    \frac{E'}{p'}
    &=
    \frac{
        \gamma'mc^2
    }{
        \gamma'mv'
    }
    \\[3pt]
    &=
    \frac{c^2}{v'}.
\end{align*}

Следовательно,
%
\begin{equation*}
    \frac{E'}{cp'}
    =
    \frac{c}{v'}.
\end{equation*}

Окончательно
%
\begin{equation*}
    \tan\theta
    =
    \frac{
        \sin\theta'
    }{
        \gamma
        \left(
            \beta\frac{c}{v'}
            +
            \cos\theta'
        \right)
    }.
\end{equation*}

Эта формула позволяет связать угол вылета частицы в движущейся
системе отсчёта с углом в лабораторной системе.

В частности, её можно использовать для задачи 1 при переходе
из системы покоя распадающейся частицы в лабораторную систему.

\textbf{Комментарий.}
В промежуточной записи появляется лишнее деление на $p'$.
Последовательное преобразование из выражений для $p_x$ и $p_y$
приводит к формуле, записанной выше.


% ============================================================
% Задача 6
% ============================================================

\newpage
\section{Задача 6. Упругое рассеяние}

Рассмотрим столкновение частицы массы $M$
с покоящейся частицей массы $m$.

\begin{rbox}[left=18pt,right=18pt,top=14pt,bottom=10pt]{[]}
    \centering

    \begin{tikzpicture}[
        >=Latex,
        every node/.style={text=rtext}
    ]
        % Лабораторная система
        \draw[
            ->,
            draw=rtext,
            line width=0.8pt
        ]
        (-0.2,0) -- (5.0,0)
        node[right] {$x$};

        \draw[
            ->,
            draw=rtext,
            line width=0.8pt
        ]
        (0,-0.2) -- (0,3.2)
        node[above] {$y$};

        % Движущаяся система
        \draw[
            ->,
            draw=rc3,
            line width=0.8pt
        ]
        (1.2,0.45) -- (5.1,0.45)
        node[right,text=rtext] {$x'$};

        \draw[
            ->,
            draw=rc3,
            line width=0.8pt
        ]
        (1.2,0.25) -- (1.2,3.1)
        node[above,text=rtext] {$y'$};

        % Скорость движущейся системы
        \draw[
            ->,
            draw=rc3,
            line width=1.0pt
        ]
        (1.8,-0.55) -- (3.6,-0.55);

        \node[
            below,
            text=rtext
        ] at (2.7,-0.55) {$\vec v$};

        % Импульс частицы
        \draw[
            ->,
            draw=rtext,
            line width=1pt
        ]
        (1.2,0.45) -- (3.7,2.5);

        \node[
            above right
        ] at (3.45,2.3) {$\vec p'$};

        % Угол
        \draw[
            draw=rtext
        ]
        (2.0,0.45)
        arc[
            start angle=0,
            end angle=39,
            radius=0.8
        ];

        \node at (2.3,0.85) {$\theta'$};

    \end{tikzpicture}

    \captionsetup{
        font=footnotesize,
        labelfont=bf,
        skip=5pt
    }

    \captionof{figure}{
        Переход от движущейся системы отсчёта
        к лабораторной.
    }
    \label{fig:practice02_angle}
\end{rbox}

Требуется найти кинетическую энергию переданной частице
как функцию угла:
%
\begin{equation*}
    T(\theta)-?
\end{equation*}

Закон сохранения четырёхимпульса:
%
\begin{equation*}
    P_0+(m,\vec 0)
    =
    P_1+P_2.
\end{equation*}

Отсюда
%
\begin{equation*}
    P_1
    =
    P_0+(m,\vec 0)-P_2.
\end{equation*}

Начальная и конечная частицы $P_0$ и $P_1$
имеют одну и ту же массу $M$, поэтому
%
\begin{equation*}
    P_0^2=P_1^2=M^2.
\end{equation*}

Для конечной частицы массы $M$
%
\begin{equation*}
    M^2
    =
    \left[
        (E_0,\vec p_0)
        +
        (m,\vec 0)
        -
        (E_2,\vec p_2)
    \right]^2.
\end{equation*}

Раскроем квадрат:
%
\begin{align*}
    M^2
    &=
    (E_0,\vec p_0)^2
    +
    (m,\vec 0)^2
    +
    (E_2,\vec p_2)^2
    \\[3pt]
    &\quad
    +
    2(E_0,\vec p_0)(m,\vec 0)
    -
    2(E_0,\vec p_0)(E_2,\vec p_2)
    -
    2(m,\vec 0)(E_2,\vec p_2).
\end{align*}

Используем
%
\begin{equation*}
    (E_0,\vec p_0)^2=M^2,
\end{equation*}
%
\begin{equation*}
    (m,\vec 0)^2=m^2,
\end{equation*}
%
и
%
\begin{equation*}
    (E_2,\vec p_2)^2=m^2.
\end{equation*}

Кроме того,
%
\begin{equation*}
    (E_0,\vec p_0)(m,\vec 0)
    =
    E_0m,
\end{equation*}
%
\begin{equation*}
    (m,\vec 0)(E_2,\vec p_2)
    =
    mE_2,
\end{equation*}
%
а
%
\begin{equation*}
    (E_0,\vec p_0)(E_2,\vec p_2)
    =
    E_0E_2-\vec p_0\cdot\vec p_2.
\end{equation*}

Если угол между $\vec p_0$ и $\vec p_2$ равен $\theta$, то
%
\begin{equation*}
    \vec p_0\cdot\vec p_2
    =
    p_0p_2\cos\theta.
\end{equation*}

Поэтому
%
\begin{equation*}
    M^2
    =
    M^2
    +
    m^2
    +
    m^2
    +
    2E_0m
    -
    2E_0E_2
    -
    2E_2m
    +
    2p_0p_2\cos\theta.
\end{equation*}

Введём кинетическую энергию частицы массы $m$:
%
\begin{equation*}
    T
    =
    E_2-m.
\end{equation*}

Её полная энергия равна
%
\begin{equation*}
    E_2
    =
    T+m.
\end{equation*}

Для начальной частицы
%
\begin{equation*}
    E_0
    =
    \sqrt{
        p_0^2+M^2
    }.
\end{equation*}

Импульс конечной частицы массы $m$:
%
\begin{equation*}
    p_2
    =
    \sqrt{
        E_2^2-m^2
    }.
\end{equation*}

Рассмотрим уравнение сохранения четырёхимпульса:

\begin{equation*}
    M^2
    =
    M^2
    +
    2m^2
    +
    2E_0m
    -
    2E_0E_2
    -
    2E_2m
    +
    2p_0p_2\cos\theta,
\end{equation*}

Сократим $M^2$, разделим уравнение сохранения на $2$
и подставим $E_2=T+m$:
%
\begin{align*}
    0
    &=m^2+E_0m-(E_0+m)(T+m)+p_0p_2\cos\theta\\
    &=-(E_0+m)T+p_0p_2\cos\theta.
\end{align*}
%
Следовательно,
%
\begin{equation*}
    (E_0+m)T=p_0p_2\cos\theta.
\end{equation*}
%
При этом
%
\begin{equation*}
    p_2^2=(T+m)^2-m^2=T(T+2m).
\end{equation*}
%
Получаем уравнение для переданной кинетической энергии:
%
\begin{equation*}
    (E_0+m)T=p_0\sqrt{T(T+2m)}\cos\theta.
\end{equation*}

При $T>0$ левая часть положительна, поэтому $\cos\theta>0$:
частица отдачи вылетает в переднюю полусферу,
$0\leq\theta<\pi/2$.
Возведём уравнение в квадрат и разделим на $T$:
%
\begin{align*}
    (E_0+m)^2T^2
    &=p_0^2T(T+2m)\cos^2\theta,\\
    \left[(E_0+m)^2-p_0^2\cos^2\theta\right]T
    &=2mp_0^2\cos^2\theta.
\end{align*}

Искомая зависимость имеет вид
%
\begin{equation*}
    \boxed{
        T(\theta)=
        \frac{2mp_0^2\cos^2\theta}
        {(E_0+m)^2-p_0^2\cos^2\theta}
    },
    \qquad E_0=\sqrt{p_0^2+M^2}.
\end{equation*}
%
Поскольку $p_0^2=E_0^2-M^2$, ответ можно записать также как
%
\begin{equation*}
    \boxed{
        T(\theta)=
        \frac{2m(E_0^2-M^2)\cos^2\theta}
        {M^2+m^2+2mE_0+(E_0^2-M^2)\sin^2\theta}
    }.
\end{equation*}

Знаменатель положителен при всех углах.
Однако значения $\theta>\pi/2$ не удовлетворяют исходному
уравнению до возведения в квадрат.
Отдельное решение $T=0$ соответствует отсутствию передачи импульса:
$p_2=0$, и угол вылета тогда не определён.
При $\theta\to\pi/2$ полученная ненулевая ветвь стремится к $T=0$.

\subsubsection*{Максимальная переданная энергия}

Энергия $T$ возрастает с $\cos^2\theta$, поэтому максимум
достигается при вылете частицы отдачи вперёд, $\theta=0$:
%
\begin{equation*}
    \boxed{
        T_{\max}=
        \frac{2mp_0^2}{M^2+m^2+2mE_0}
    }.
\end{equation*}
%
Если начальная кинетическая энергия налетающей частицы
$T_0=E_0-M$, то
%
\begin{equation*}
    T_{\max}=
    \frac{2mT_0(T_0+2M)}{(M+m)^2+2mT_0}.
\end{equation*}
%
В частности, при $M=m$ получаем $T_{\max}=T_0$:
возможна передача всей начальной кинетической энергии.

\subsubsection*{Проверка нерелятивистского предела}

При $T_0\ll M$ имеем $p_0^2\simeq2MT_0$ и $E_0\simeq M$.
Сохраняя главный порядок по $T_0/M$, находим
%
\begin{equation*}
    T(\theta)\simeq
    \frac{4mM}{(M+m)^2}\,T_0\cos^2\theta,
    \qquad
    T_{\max}\simeq\frac{4mM}{(M+m)^2}\,T_0.
\end{equation*}
%
Это совпадает с результатом нерелятивистской кинематики
упругого столкновения.

\end{document}
